%% LyX 2.4.4 created this file.  For more info, see https://www.lyx.org/.
%% Do not edit unless you really know what you are doing.
\documentclass[10pt,english]{beamer}
\usepackage{mathptmx}
\usepackage[T1]{fontenc}
\usepackage[latin9]{inputenc}
\usepackage{units}
\usepackage{amssymb}
\usepackage{graphicx}

\makeatletter

%%%%%%%%%%%%%%%%%%%%%%%%%%%%%% LyX specific LaTeX commands.
%% A simple dot to overcome graphicx limitations
\newcommand{\lyxdot}{.}


%%%%%%%%%%%%%%%%%%%%%%%%%%%%%% Textclass specific LaTeX commands.
% this default might be overridden by plain title style
\newcommand\makebeamertitle{\frame{\maketitle}}%
% (ERT) argument for the TOC
\AtBeginDocument{%
  \let\origtableofcontents=\tableofcontents
  \def\tableofcontents{\@ifnextchar[{\origtableofcontents}{\gobbletableofcontents}}
  \def\gobbletableofcontents#1{\origtableofcontents}
}

%%%%%%%%%%%%%%%%%%%%%%%%%%%%%% User specified LaTeX commands.
\usetheme{Warsaw}
% or ...

\usepackage{pgfpages,pgfnodes,pgfarrows,pgfheaps}

\setbeamercovered{transparent}
% or whatever (possibly just delete it)
\setbeamercovered{invisible} %makes overlays invisible


%gets rid of navigation symbols
\setbeamertemplate{navigation symbols}{}
\usepackage{hyperref}

%%%setshandout pages per page and adds frames
%\usepackage{pgfpages}
%\pgfpagesuselayout{4 on 1}[a4paper, landscape, border shrink=7mm]
%\pgfpageslogicalpageoptions{1}{border code=\pgfusepath{stroke}}
%\pgfpageslogicalpageoptions{2}{border code=\pgfusepath{stroke}}
%\pgfpageslogicalpageoptions{3}{border code=\pgfusepath{stroke}}
%\pgfpageslogicalpageoptions{4}{border code=\pgfusepath{stroke}}

\makeatother

\usepackage{babel}
\begin{document}
%sets caption options. Requires \usepackage{caption} in preamble.
%\captionsetup[figure]{labelformat=empty, labelsep=none}
%the default beamer dimensions are 12.8 x 9.6cm.
\title{Generalized Branching in Circle Packing}
\subtitle{A Discrete Ahflors Function}
\author{Edward Crane, Ken Stephenson, and James Ashe}
\titlegraphic{\includegraphics[height=1.5cm]{D:/home/jim/JamesAshe/JCSU/images/jcsu_bull.pdf}}
\institute{Johnson C. Smith University}
\date{{}}

\makebeamertitle

\beamerdefaultoverlayspecification{<+->}
\begin{frame}{\textbf{What is a Circle Packing?}}

\uncover<2->{\textbf{Geometry + Complex = Circle Packing}}
\begin{columns}


\column{6.75cm }

\includegraphics[width=6.75cm]{../Thesis_defense/pictures/owl_1}<1|handout:0>\includegraphics[width=6.75cm]{pics/AMS_Knoxville_2014_pics/owl_4_colord_border_circles}<2|handout:0>\includegraphics[width=6.75cm]{pics/AMS_Knoxville_2014_pics/owl_max_colord_border_circles}<3|handout:0>\includegraphics[width=6.75cm]{pics/AMS_Knoxville_2014_pics/owl_4_sphere}<4|handout:0>
\end{columns}
\end{frame}

\begin{frame}{Conformal Mapping}

\begin{center}
\includegraphics[width=10cm]{D:/home/jim/JamesAshe/pictures/Riemann_Map.jpg}
\par\end{center}

\uncover<1->{When simply connected.. everything works beautifully.}

\uncover<2->{When not simply connected the current theory does not
suffice} 
\end{frame}

\begin{frame}{Local coherence and branching}

\begin{columns}


\column{6cm}
\begin{exampleblock}<1->{Flowers with angle sums of}

\begin{itemize}
\item <2-8>$2\pi$, a (trivially) branched flower.
\item <9-15>$4\pi$, a branched flower.
\item <16->$2.75\pi$, an incoherent flower
\end{itemize}
\end{exampleblock}

\column{5cm}

\includegraphics[width=5cm]{../Thesis_defense/aux_pics/cp_5587/cpShot-046}<2|handout:0>\includegraphics[width=5cm]{../Thesis_defense/aux_pics/cp_5587/cpShot-047}<3|handout:0>\includegraphics[width=5cm]{../Thesis_defense/aux_pics/cp_5587/cpShot-048}<4|handout:0>\includegraphics[width=5cm]{../Thesis_defense/aux_pics/cp_5587/cpShot-049}<5|handout:0>\includegraphics[width=5cm]{../Thesis_defense/aux_pics/cp_5587/cpShot-050}<6|handout:0>\includegraphics[width=5cm]{../Thesis_defense/aux_pics/cp_5587/cpShot-051}<7|handout:0>\includegraphics[width=5cm]{../Thesis_defense/aux_pics/cp_5587/cpShot-052}<8|handout:0>\includegraphics[width=5cm]{../Thesis_defense/aux_pics/cp_5587/cpShot4pi-053}<9|handout:0>\includegraphics[width=5cm]{../Thesis_defense/aux_pics/cp_5587/cpShot4pi-054}<10|handout:0>\includegraphics[width=5cm]{../Thesis_defense/aux_pics/cp_5587/cpShot4pi-055}<11|handout:0>\includegraphics[width=5cm]{../Thesis_defense/aux_pics/cp_5587/cpShot4pi-056}<12|handout:0>\includegraphics[width=5cm]{../Thesis_defense/aux_pics/cp_5587/cpShot4pi-057}<13|handout:0>\includegraphics[width=5cm]{../Thesis_defense/aux_pics/cp_5587/cpShot4pi-058}<14|handout:0>\includegraphics[width=5cm]{../Thesis_defense/aux_pics/cp_5587/cpShot4pi-059}<15|handout:0>\includegraphics[width=5cm]{../Thesis_defense/aux_pics/cp_5587/cpShot2-75pi-067}<16|handout:0>\includegraphics[width=5cm]{../Thesis_defense/aux_pics/cp_5587/cpShot2-75pi-068}<17|handout:0>\includegraphics[width=5cm]{../Thesis_defense/aux_pics/cp_5587/cpShot2-75pi-069}<18|handout:0>\includegraphics[width=5cm]{../Thesis_defense/aux_pics/cp_5587/cpShot2-75pi-070}<19|handout:0>\includegraphics[width=5cm]{../Thesis_defense/aux_pics/cp_5587/cpShot2-75pi-071}<20|handout:0>\includegraphics[width=5cm]{../Thesis_defense/aux_pics/cp_5587/cpShot2-75pi-072}<21|handout:0>\includegraphics[width=5cm]{../Thesis_defense/aux_pics/cp_5587/cpShot2-75pi-073}<22|handout:0>\includegraphics[width=5cm]{../Thesis_defense/aux_pics/cp_5587/cpShot2-75pi-074}<23|handout:0>\includegraphics[width=5cm]{../Thesis_defense/aux_pics/cp_5587/cpShot2-75pi-075}<24->
\end{columns}
\end{frame}

\begin{frame}{Local coherence $\nRightarrow$ coherence}

\begin{columns}


\column{6.5cm}

\includegraphics[height=6cm]{../Thesis_defense/pictures/cp_5587/cpShot2-75pi-080}<1|handout:0>\includegraphics[height=6cm]{../Thesis_defense/pictures/cp_5587/cpShot2-75pi-081}<2|handout:0>\includegraphics[height=6cm]{../Thesis_defense/pictures/cp_5587/cpShot2-75pi-082}<3|handout:0>\includegraphics[height=6cm]{../Thesis_defense/pictures/cp_5587/cpShot2-75pi-083}<4|handout:0>\includegraphics[height=6cm]{../Thesis_defense/pictures/cp_5587/cpShot2-75pi-084}<5|handout:0>\includegraphics[height=6cm]{../Thesis_defense/pictures/cp_5587/cpShot2-75pi-085}<6|handout:0>

\column{4.5cm}
\begin{exampleblock}<1-2>{Example}

A circle packing.
\end{exampleblock}
{}
\begin{exampleblock}<3-6>{Example}

Not a circle packing.
\end{exampleblock}
\end{columns}
\end{frame}

\begin{frame}{Discretization issues}

\begin{block}<only@1-20|visible@1->{\textbf{Ahlfors function:} maximal packing of an annulus.}

\begin{itemize}
\item <1>What we want.
\item <2-20>What we get after symmetry is lost.
\end{itemize}
\end{block}
\begin{columns}


\column{8.5cm}
\begin{center}
\includegraphics[width=9.5cm]{../Thesis_defense/pictures/cp_19687/cpShot-001}<1|handout:0>\includegraphics[width=9.5cm]{../Thesis_defense/pictures/cp_19687/cpShot-002}<2|handout:0>\includegraphics[height=6cm]{../Thesis_defense/pictures/cp_3925/cpShot-001}<3|handout:1>\includegraphics[height=6cm]{../Thesis_defense/pictures/cp_3925/cpShot-025}<4|handout:0>\includegraphics[height=6cm]{../Thesis_defense/pictures/cp_3925/cpShot-026}<5|handout:0>\includegraphics[height=6cm]{../Thesis_defense/pictures/cp_3925/cpShot-027}<6|handout:0>\includegraphics[height=6cm]{../Thesis_defense/pictures/cp_3925/cpShot-028}<7|handout:0>\includegraphics[height=6cm]{../Thesis_defense/pictures/cp_3925/cpShot-029}<8|handout:0>\includegraphics[height=6cm]{../Thesis_defense/pictures/cp_3925/cpShot-030}<9|handout:0>\includegraphics[height=6cm]{../Thesis_defense/pictures/cp_3925/cpShot-031}<10|handout:0>\includegraphics[height=6cm]{../Thesis_defense/pictures/cp_3925/cpShot-032}<11|handout:0>\includegraphics[height=6cm]{../Thesis_defense/pictures/cp_3925/cpShot-033}<12|handout:0>\includegraphics[height=6cm]{../Thesis_defense/pictures/cp_3925/cpShot-034}<13|handout:0>\includegraphics[height=6cm]{../Thesis_defense/pictures/cp_3925/cpShot-035}<14|handout:0>\includegraphics[height=6cm]{../Thesis_defense/pictures/cp_3925/cpShot-036}<15|handout:0>\includegraphics[height=6cm]{../Thesis_defense/pictures/cp_3925/cpShot-039}<16|handout:0>\includegraphics[height=6cm]{../Thesis_defense/pictures/cp_3925/cpShot-040}<17|handout:0>\includegraphics[height=6cm]{../Thesis_defense/pictures/cp_3925/cpShot-041}<18|handout:0>\includegraphics[height=6cm]{../Thesis_defense/pictures/cp_3925/cpShot-042}<19|handout:0>\includegraphics[height=6cm]{../Thesis_defense/pictures/cp_3925/cpShot-043}<20|handout:2>
\par\end{center}

\end{columns}
\end{frame}

\begin{frame}{Generalized branching}

\begin{block}<only@1-6|visible@1->{Definition}

A simply connected subcomplex of $K$ (with decorations) is a \emph{generalized
branch point} for a label $R$ if its border chain has trivial holonomy.
\end{block}
\begin{columns}


\column{5cm}
\begin{itemize}
\item <2->The idea is to allow areas of ambiguity so that flexibility can
be introduced.
\item <3->The ambiguity is contained within the subcomplex.
\end{itemize}

\column{5cm}

\includegraphics[width=5cm]{D:/home/jim/JamesAshe/talks/AMS_sectional_Knoxville_2014/thesis_pictures/gen_branch_blank.svg}<1>\includegraphics[width=5cm]{D:/home/jim/JamesAshe/talks/AMS_sectional_Knoxville_2014/thesis_pictures/ambiguity.svg}<2>\includegraphics[width=5cm]{D:/home/jim/JamesAshe/talks/AMS_sectional_Knoxville_2014/thesis_pictures/exterior_face_with_branching_c.svg}<3>\includegraphics[width=5cm]{D:/home/jim/JamesAshe/talks/AMS_sectional_Knoxville_2014/thesis_pictures/exterior_face_with_branching_d.svg}<4->
\end{columns}
\end{frame}

\begin{frame}{Generalized branching}

\begin{columns}


\column{6.0cm}
\begin{block}{A decorated complex $K(I,A)$}

\begin{description}
\item [{Complex~$K$:}] <1->an abstract simplicial $2$-complex which
triangulates an oriented topological surface.
\item [{Angle~Assignment~$A$:}] <2->a collection of angle values paired
with each interior vertex of $K$.
\item [{Edge~Assignment~$I$:}] <3->a collection of M�bius invariant
inversive distances paired with each edge in $K$.
\end{description}
\end{block}

\column{5.0cm}
\begin{overlayarea}{}{}
{5.0cm}{6.5cm}

\only<1>{
\begin{figure}
\begin{centering}
\includegraphics[width=5cm]{D:/home/jim/JamesAshe/talks/AMS_sectional_Knoxville_2014/thesis_pictures/basic_K_1.svg}
\par\end{centering}
\caption{$K$ tells of us which circles are neighbors.}
\end{figure}
}

\only<2>{
\begin{figure}
\begin{centering}
\includegraphics[width=5cm]{D:/home/jim/JamesAshe/talks/AMS_sectional_Knoxville_2014/thesis_pictures/basic_A_3.svg}
\par\end{centering}
\caption{$A$ tells us what the resulting angles should be.}
\end{figure}
}

\only<3>{
\begin{figure}
\begin{centering}
\includegraphics[width=5cm]{D:/home/jim/JamesAshe/talks/AMS_sectional_Knoxville_2014/thesis_pictures/basic_I_2.svg}
\par\end{centering}
\caption{$I$ tells of us how far apart circles are supposed to be placed.}
\end{figure}
}
\end{overlayarea}

\end{columns}
\end{frame}

\begin{frame}{Singular Branching (interstice of face)}

\begin{columns}


\column{6.5cm}

\includegraphics[width=6.5cm]{\string"pics/AMS_Knoxville_2014_pics/singular_example_center-interstice/sing 1\string".eps}<1|handout:0>\includegraphics[width=6.5cm]{\string"pics/AMS_Knoxville_2014_pics/singular_example_center-interstice/sing 2\string".eps}<2|handout:0>\includegraphics[width=6.5cm]{\string"pics/AMS_Knoxville_2014_pics/singular_example_center-interstice/sing 3\string".eps}<3|handout:0>\includegraphics[width=6.5cm]{\string"pics/AMS_Knoxville_2014_pics/singular_example_center-interstice/sing 4\string".eps}<4|handout:0>\includegraphics[width=6.5cm]{\string"pics/AMS_Knoxville_2014_pics/singular_example_center-interstice/sing 5\string".eps}<5|handout:0>\includegraphics[width=6.5cm]{\string"pics/AMS_Knoxville_2014_pics/singular_example_center-interstice/sing 6\string".eps}<6|handout:0>\includegraphics[width=6.5cm]{\string"pics/AMS_Knoxville_2014_pics/singular_example_center-interstice/sing 7\string".eps}<7|handout:0>\includegraphics[width=6.5cm]{\string"pics/AMS_Knoxville_2014_pics/singular_example_center-interstice/sing 8\string".eps}<8|handout:0>\includegraphics[width=6.5cm]{\string"pics/AMS_Knoxville_2014_pics/singular_example_center-interstice/sing 9\string".eps}<9|handout:0>\includegraphics[width=6.5cm]{\string"pics/AMS_Knoxville_2014_pics/singular_example_center-interstice/sing 10\string".eps}<10|handout:0>\includegraphics[width=6.5cm]{\string"pics/AMS_Knoxville_2014_pics/singular_example_center-interstice/sing 11\string".eps}<11|handout:0>\includegraphics[width=6.5cm]{\string"pics/AMS_Knoxville_2014_pics/singular_example_center-interstice/sing 12\string".eps}<12|handout:0>\includegraphics[width=6.5cm]{\string"pics/AMS_Knoxville_2014_pics/singular_example_center-interstice/sing 13\string".eps}<13|handout:0>\includegraphics[width=6.5cm]{\string"pics/AMS_Knoxville_2014_pics/singular_example_center-interstice/sing 14\string".eps}<14|handout:0>\includegraphics[width=6.5cm]{\string"pics/AMS_Knoxville_2014_pics/singular_example_center-interstice/sing 15\string".eps}<15|handout:0>\includegraphics[width=6.5cm]{pics/AMS_Knoxville_2014_pics/singular_example_center-interstice/sing16}<16|handout:0>\includegraphics[width=6.5cm]{pics/AMS_Knoxville_2014_pics/singular_example_center-interstice/sing17}<17|handout:0>\includegraphics[width=6.5cm]{pics/AMS_Knoxville_2014_pics/singular_example_center-interstice/sing18}<18|handout:0>\includegraphics[width=6.5cm]{pics/AMS_Knoxville_2014_pics/singular_example_center-interstice/sing19}<19|handout:0>\includegraphics[width=6.5cm]{pics/AMS_Knoxville_2014_pics/singular_example_center-interstice/sing20}<20|handout:0>\includegraphics[width=6.5cm]{pics/AMS_Knoxville_2014_pics/singular_example_center-interstice/sing21}<21|handout:0>\includegraphics[width=6.5cm]{pics/AMS_Knoxville_2014_pics/singular_example_center-interstice/sing22}<22|handout:0>\includegraphics[width=6.5cm]{pics/AMS_Knoxville_2014_pics/singular_example_center-interstice/sing23}<23|handout:0>\includegraphics[width=6.5cm]{pics/AMS_Knoxville_2014_pics/singular_example_center-interstice/sing24}<24|handout:0>\includegraphics[width=6.5cm]{pics/AMS_Knoxville_2014_pics/singular_example_center-interstice/sing25}<25|handout:0>\includegraphics[width=6.5cm]{pics/AMS_Knoxville_2014_pics/singular_example_center-interstice/sing26}<26|handout:0>\includegraphics[width=6.5cm]{pics/AMS_Knoxville_2014_pics/singular_example_center-interstice/sing27}<27|handout:0>\includegraphics[width=6.5cm]{pics/AMS_Knoxville_2014_pics/singular_example_center-interstice/sing28}<28|handout:0>\includegraphics[width=6.5cm]{pics/AMS_Knoxville_2014_pics/singular_example_center-interstice/sing30}<29|handout:0>

\column{4.5cm}
\begin{block}{}

\begin{itemize}
\item <1->Auxiliary circles are added, including one at the ``branch point''. 
\item <1->''branch circle'' radius is forced to $0$.
\item <1->Resulting overlaps provide our parameters. 
\end{itemize}
\end{block}
\end{columns}
\end{frame}

\begin{frame}{Chaperone Branching (in a circle)}

\begin{columns}


\column{6.5cm}

\includegraphics[width=6.5cm]{pics/AMS_Knoxville_2014_pics/Chaperone_labeled_just_hex}<1|handout:0>\includegraphics[width=6.5cm]{pics/AMS_Knoxville_2014_pics/Chaperone_labeled}<2|handout:0>\includegraphics[width=6.5cm]{pics/AMS_Knoxville_2014_pics/chap_unbranched_movie/cpShot-154}<3|handout:0>\includegraphics[width=6.5cm]{pics/AMS_Knoxville_2014_pics/chap_unbranched_movie/cpShot-156}<4|handout:0>\includegraphics[width=6.5cm]{pics/AMS_Knoxville_2014_pics/chap_unbranched_movie/cpShot-157}<5|handout:0>\includegraphics[width=6.5cm]{pics/AMS_Knoxville_2014_pics/chap_unbranched_movie/cpShot-158}<6|handout:0>\includegraphics[width=6.5cm]{pics/AMS_Knoxville_2014_pics/chap_unbranched_movie/cpShot-159}<7|handout:0>\includegraphics[width=6.5cm]{pics/AMS_Knoxville_2014_pics/chap_unbranched_movie/cpShot-160}<8|handout:0>\includegraphics[width=6.5cm]{pics/AMS_Knoxville_2014_pics/chap_unbranched_movie/cpShot-161}<9|handout:0>\includegraphics[width=6.5cm]{pics/AMS_Knoxville_2014_pics/chap_unbranched_movie/cpShot-162}<10|handout:0>\includegraphics[width=6.5cm]{pics/AMS_Knoxville_2014_pics/chap_unbranched_movie/cpShot-163}<11|handout:0>\includegraphics[width=6.5cm]{pics/AMS_Knoxville_2014_pics/chap_unbranched_movie/cpShot-164}<12|handout:0>\includegraphics[width=6.5cm]{pics/AMS_Knoxville_2014_pics/chap_unbranched_movie/cpShot-165}<13|handout:0>\includegraphics[width=6.5cm]{pics/AMS_Knoxville_2014_pics/chap_unbranched_movie/cpShot-167}<14|handout:0>\includegraphics[width=6.5cm]{pics/AMS_Knoxville_2014_pics/chap_unbranched_movie/cpShot-168}<15|handout:0>\includegraphics[width=6.5cm]{pics/AMS_Knoxville_2014_pics/chap_unbranched_movie/cpShot-169}<16|handout:0>\includegraphics[width=6.5cm]{pics/AMS_Knoxville_2014_pics/chap_unbranched_movie/cpShot-170}<17|handout:0>\includegraphics[width=6.5cm]{pics/AMS_Knoxville_2014_pics/chap_unbranched_movie/cpShot-171}<18|handout:0>\includegraphics[width=6.5cm]{pics/AMS_Knoxville_2014_pics/chap_unbranched_movie/cpShot-172}<19|handout:0>\includegraphics[width=6.5cm]{pics/AMS_Knoxville_2014_pics/chap_unbranched_movie/cpShot-173}<20|handout:0>\includegraphics[width=6.5cm]{pics/AMS_Knoxville_2014_pics/chap_unbranched_movie/cpShot-174}<21|handout:0>\includegraphics[width=6.5cm]{pics/AMS_Knoxville_2014_pics/chap_unbranched_movie/cpShot-175}<22|handout:0>

\column{4.5cm}
\begin{block}{}

\begin{itemize}
\item <1->Auxiliary circles are introduced, one at the ``branch point''. 
\item <1->Overlap of $2\pi$ is distributed on each Chaperone circle.
\end{itemize}
\end{block}
\end{columns}
\end{frame}

\begin{frame}{Chaperone Branching}

\begin{columns}


\column{6.5cm}

\includegraphics[width=6.5cm]{pics/AMS_Knoxville_2014_pics/chap_branched_movie/cpShot-220}<1|handout:0>\includegraphics[width=6.5cm]{pics/AMS_Knoxville_2014_pics/chap_branched_movie/cpShot-221}<2|handout:0>\includegraphics[width=6.5cm]{pics/AMS_Knoxville_2014_pics/chap_branched_movie/cpShot-222}<3|handout:0>\includegraphics[width=6.5cm]{pics/AMS_Knoxville_2014_pics/chap_branched_movie/cpShot-223}<4|handout:0>\includegraphics[width=6.5cm]{pics/AMS_Knoxville_2014_pics/chap_branched_movie/cpShot-224}<5|handout:0>\includegraphics[width=6.5cm]{pics/AMS_Knoxville_2014_pics/chap_branched_movie/cpShot-225}<6|handout:0>\includegraphics[width=6.5cm]{pics/AMS_Knoxville_2014_pics/chap_branched_movie/cpShot-226}<7|handout:0>\includegraphics[width=6.5cm]{pics/AMS_Knoxville_2014_pics/chap_branched_movie/cpShot-227}<8|handout:0>\includegraphics[width=6.5cm]{pics/AMS_Knoxville_2014_pics/chap_branched_movie/cpShot-228}<9|handout:0>\includegraphics[width=6.5cm]{pics/AMS_Knoxville_2014_pics/chap_branched_movie/cpShot-229}<10|handout:0>\includegraphics[width=6.5cm]{pics/AMS_Knoxville_2014_pics/chap_branched_movie/cpShot-231}<11|handout:0>\includegraphics[width=6.5cm]{pics/AMS_Knoxville_2014_pics/chap_branched_movie/cpShot-232}<12|handout:0>\includegraphics[width=6.5cm]{pics/AMS_Knoxville_2014_pics/chap_branched_movie/cpShot-233}<13|handout:0>\includegraphics[width=6.5cm]{pics/AMS_Knoxville_2014_pics/chap_branched_movie/cpShot-234}<14|handout:0>\includegraphics[width=6.5cm]{pics/AMS_Knoxville_2014_pics/chap_branched_movie/cpShot-235}<15|handout:0>\includegraphics[width=6.5cm]{pics/AMS_Knoxville_2014_pics/chap_branched_movie/cpShot-236}<16|handout:0>\includegraphics[width=6.5cm]{pics/AMS_Knoxville_2014_pics/chap_branched_movie/cpShot-237}<17|handout:0>\includegraphics[width=6.5cm]{pics/AMS_Knoxville_2014_pics/chap_branched_movie/cpShot-238}<18|handout:0>\includegraphics[width=6.5cm]{pics/AMS_Knoxville_2014_pics/chap_branched_movie/cpShot-239}<19|handout:0>\includegraphics[width=6.5cm]{pics/AMS_Knoxville_2014_pics/chap_branched_movie/cpShot-240}<20|handout:0>\includegraphics[width=6.5cm]{pics/AMS_Knoxville_2014_pics/chap_branched_movie/cpShot-241}<21|handout:0>\includegraphics[width=6.5cm]{pics/AMS_Knoxville_2014_pics/chap_branched_movie/cpShot-244}<22|handout:0>

\column{4.5cm}
\begin{block}{}

\begin{itemize}
\item <1->Auxiliary circles are introduced, one at the ``branch point''. 
\item <1->Overlap of $2\pi$ is distributed on each Chaperone circle.
\end{itemize}
\end{block}
\end{columns}
\end{frame}

\begin{frame}{Triangles with deep overlaps}

\begin{columns}


\column{5.5cm}
\begin{block}{Monotonicity Lemma}
If the edge assignments of a face are 

\begin{itemize}
\item []in $[0,\nicefrac{\pi}{2}]$
\item []OR sum to $\leq\pi$,
\end{itemize}
then for any three radii in $\mathbb{C}$ or $\mathbb{D}$:
\begin{enumerate}
\item <2->[1.]A triangle with these overlaps exists.
\item <3->[2.]The triangle is monotone.
\end{enumerate}
\end{block}
\begin{theorem}
9->{}The radii for a CP with singular/chaperone branching exist
(under some restrictions) and are unique.
\end{theorem}


\column{6.0cm}

\noindent\begin{minipage}[t][7cm]{1\columnwidth}%
\begin{center}
\includegraphics[height=4.5cm]{D:/home/jim/JamesAshe/talks/AMS_sectional_Knoxville_2014/thesis_pictures/cp_8207/cpShot-009.jpg}<1-3>\includegraphics[height=4.5cm]{D:/home/jim/JamesAshe/talks/AMS_sectional_Knoxville_2014/thesis_pictures/cp_8207/cpShot-009.jpg}<4>\includegraphics[height=4.5cm]{D:/home/jim/JamesAshe/talks/AMS_sectional_Knoxville_2014/thesis_pictures/cp_8207/cpShot-010.jpg}<5>\includegraphics[height=4.5cm]{D:/home/jim/JamesAshe/talks/AMS_sectional_Knoxville_2014/thesis_pictures/cp_8207/cpShot-011.jpg}<6>\includegraphics[height=4.5cm]{D:/home/jim/JamesAshe/talks/AMS_sectional_Knoxville_2014/thesis_pictures/cp_8207/cpShot-012.jpg}<7>\includegraphics[height=4.5cm]{D:/home/jim/JamesAshe/talks/AMS_sectional_Knoxville_2014/thesis_pictures/cp_8207/cpShot-013.jpg}<8-9>
\par\end{center}%
\end{minipage}
\end{columns}
\end{frame}

\begin{frame}{Continuously moving a branch point}

\uncover<1>{Traditional branching at vertex}

\uncover<22>{Traditional branching at vertex}
\begin{columns}


\column{5.5cm}

\includegraphics[width=5.5cm]{pics/AMS_Knoxville_2014_pics/gen_branch_movie/cpShot-082}<1|handout:0>\includegraphics[width=5.5cm]{pics/AMS_Knoxville_2014_pics/gen_branch_movie/cpShot-085}<2|handout:0>\includegraphics[width=5.5cm]{pics/AMS_Knoxville_2014_pics/gen_branch_movie/cpShot-086}<3|handout:0>\includegraphics[width=5.5cm]{pics/AMS_Knoxville_2014_pics/gen_branch_movie/cpShot-087}<4|handout:0>\includegraphics[width=5.5cm]{pics/AMS_Knoxville_2014_pics/gen_branch_movie/cpShot-088}<5|handout:0>\includegraphics[width=5.5cm]{pics/AMS_Knoxville_2014_pics/gen_branch_movie/cpShot-090}<6|handout:0>\includegraphics[width=5.5cm]{pics/AMS_Knoxville_2014_pics/gen_branch_movie/cpShot-091}<7|handout:0>\includegraphics[width=5.5cm]{pics/AMS_Knoxville_2014_pics/gen_branch_movie/cpShot-092}<8|handout:0>\includegraphics[width=5.5cm]{pics/AMS_Knoxville_2014_pics/gen_branch_movie/cpShot-093}<9|handout:0>\includegraphics[width=5.5cm]{pics/AMS_Knoxville_2014_pics/gen_branch_movie/cpShot-094}<10|handout:0>\includegraphics[width=5.5cm]{pics/AMS_Knoxville_2014_pics/gen_branch_movie/cpShot-095}<11|handout:0>\includegraphics[width=5.5cm]{pics/AMS_Knoxville_2014_pics/gen_branch_movie/cpShot-096}<12|handout:0>\includegraphics[width=5.5cm]{pics/AMS_Knoxville_2014_pics/gen_branch_movie/cpShot-098}<13|handout:0>\includegraphics[width=5.5cm]{pics/AMS_Knoxville_2014_pics/gen_branch_movie/cpShot-099}<14|handout:0>\includegraphics[width=5.5cm]{pics/AMS_Knoxville_2014_pics/gen_branch_movie/cpShot-100}<15|handout:0>\includegraphics[width=5.5cm]{pics/AMS_Knoxville_2014_pics/gen_branch_movie/cpShot-101}<16|handout:0>\includegraphics[width=5.5cm]{pics/AMS_Knoxville_2014_pics/gen_branch_movie/cpShot-102}<17|handout:0>\includegraphics[width=5.5cm]{pics/AMS_Knoxville_2014_pics/gen_branch_movie/cpShot-103}<18|handout:0>\includegraphics[width=5.5cm]{pics/AMS_Knoxville_2014_pics/gen_branch_movie/cpShot-104}<19|handout:0>\includegraphics[width=5.5cm]{pics/AMS_Knoxville_2014_pics/gen_branch_movie/cpShot-105}<20|handout:0>\includegraphics[width=5.5cm]{pics/AMS_Knoxville_2014_pics/gen_branch_movie/cpShot-106}<21|handout:0>\includegraphics[width=5.5cm]{pics/AMS_Knoxville_2014_pics/gen_branch_movie/cpShot-107}<22|handout:0>

\column{5.5cm}

\includegraphics[width=5.5cm]{pics/AMS_Knoxville_2014_pics/gen_branch_movie/domain_path\lyxdot jpg}
\end{columns}
\end{frame}

\begin{frame}{A Discrete Ahlfors Function}

\begin{columns}


\column{5cm}
\begin{exampleblock}<1-14>{Example}

An annulus with no translational symmetry.
\end{exampleblock}
{}
\begin{exampleblock}<15-25>{Example}

A non-symmetric annulus.
\end{exampleblock}

\column{7cm}

\noindent\begin{minipage}[c][6.5cm]{1\columnwidth}%
\begin{flushleft}
\includegraphics[width=5.5cm]{D:/home/jim/JamesAshe/talks/AMS_sectional_Knoxville_2014/thesis_pictures/cp_28133/cpShot-000.jpg}<1>\includegraphics[width=5.5cm]{D:/home/jim/JamesAshe/talks/AMS_sectional_Knoxville_2014/thesis_pictures/cp_28133/cpShot-001.jpg}<2>\includegraphics[width=5.5cm]{D:/home/jim/JamesAshe/talks/AMS_sectional_Knoxville_2014/thesis_pictures/cp_28133/cpShot-002.jpg}<3>\includegraphics[width=5.5cm]{D:/home/jim/JamesAshe/talks/AMS_sectional_Knoxville_2014/thesis_pictures/cp_28133/cpShot-003.jpg}<4-5>\includegraphics[width=5.5cm]{D:/home/jim/JamesAshe/talks/AMS_sectional_Knoxville_2014/thesis_pictures/cp_28133/cpShot-006.jpg}<6>\includegraphics[width=5.5cm]{D:/home/jim/JamesAshe/talks/AMS_sectional_Knoxville_2014/thesis_pictures/cp_28133/cpShot-007.jpg}<7>\includegraphics[width=5.5cm]{D:/home/jim/JamesAshe/talks/AMS_sectional_Knoxville_2014/thesis_pictures/cp_28133/cpShot-008.jpg}<8>\includegraphics[width=5.5cm]{D:/home/jim/JamesAshe/talks/AMS_sectional_Knoxville_2014/thesis_pictures/cp_28133/cpShot-009.jpg}<9>\includegraphics[width=5.5cm]{D:/home/jim/JamesAshe/talks/AMS_sectional_Knoxville_2014/thesis_pictures/cp_28133/cpShot-010.jpg}<10>\includegraphics[width=5.5cm]{D:/home/jim/JamesAshe/talks/AMS_sectional_Knoxville_2014/thesis_pictures/cp_28133/cpShot-011.jpg}<11>\includegraphics[width=5.5cm]{D:/home/jim/JamesAshe/talks/AMS_sectional_Knoxville_2014/thesis_pictures/cp_28133/cpShot-012.jpg}<12>\includegraphics[width=5.5cm]{D:/home/jim/JamesAshe/talks/AMS_sectional_Knoxville_2014/thesis_pictures/cp_28133/cpShot-013.jpg}<13>\includegraphics[width=5.5cm]{D:/home/jim/JamesAshe/talks/AMS_sectional_Knoxville_2014/thesis_pictures/cp_28133/cpShot-014.jpg}<14>\includegraphics[width=5.5cm]{D:/home/jim/JamesAshe/talks/AMS_sectional_Knoxville_2014/thesis_pictures/cp_28134/cpShot-015.jpg}<15>\includegraphics[width=5.5cm]{D:/home/jim/JamesAshe/talks/AMS_sectional_Knoxville_2014/thesis_pictures/cp_28134/cpShot-016.jpg}<16>\includegraphics[width=5.5cm]{D:/home/jim/JamesAshe/talks/AMS_sectional_Knoxville_2014/thesis_pictures/cp_28134/cpShot-019.jpg}<17>\includegraphics[width=5.5cm]{D:/home/jim/JamesAshe/talks/AMS_sectional_Knoxville_2014/thesis_pictures/cp_28134/cpShot-020.jpg}<18>\includegraphics[width=5.5cm]{D:/home/jim/JamesAshe/talks/AMS_sectional_Knoxville_2014/thesis_pictures/cp_28134/cpShot-021.jpg}<19>\includegraphics[width=5.5cm]{D:/home/jim/JamesAshe/talks/AMS_sectional_Knoxville_2014/thesis_pictures/cp_28134/cpShot-022.jpg}<20>\includegraphics[width=5.5cm]{D:/home/jim/JamesAshe/talks/AMS_sectional_Knoxville_2014/thesis_pictures/cp_28134/cpShot-023.jpg}<21>\includegraphics[width=5.5cm]{D:/home/jim/JamesAshe/talks/AMS_sectional_Knoxville_2014/thesis_pictures/cp_28134/cpShot-024.jpg}<22>\includegraphics[width=5.5cm]{D:/home/jim/JamesAshe/talks/AMS_sectional_Knoxville_2014/thesis_pictures/cp_28134/cpShot-025.jpg}<23>\includegraphics[width=5.5cm]{D:/home/jim/JamesAshe/talks/AMS_sectional_Knoxville_2014/thesis_pictures/cp_28134/cpShot-026.jpg}<24>\includegraphics[width=5.5cm]{D:/home/jim/JamesAshe/talks/AMS_sectional_Knoxville_2014/thesis_pictures/cp_28134/cpShot-029.jpg}<25>
\par\end{flushleft}%
\end{minipage}
\end{columns}
\end{frame}

\begin{frame}{Another Application}

\begin{columns}


\column{5cm}
\begin{exampleblock}<1-15>{Example}

A quadratic polynomial.
\end{exampleblock}

\column{7cm}

\noindent\begin{minipage}[c][6.5cm]{1\columnwidth}%
\begin{flushleft}
\includegraphics[width=5.5cm]{D:/home/jim/JamesAshe/talks/AMS_sectional_Knoxville_2014/thesis_pictures/cp_28501/cpShot-000.jpg}<1-2>\includegraphics[width=5.5cm]{D:/home/jim/JamesAshe/talks/AMS_sectional_Knoxville_2014/thesis_pictures/cp_28501/cpShot-001.jpg}<3>\includegraphics[width=5.5cm]{D:/home/jim/JamesAshe/talks/AMS_sectional_Knoxville_2014/thesis_pictures/cp_28501/cpShot-002.jpg}<4>\includegraphics[width=5.5cm]{D:/home/jim/JamesAshe/talks/AMS_sectional_Knoxville_2014/thesis_pictures/cp_28501/cpShot-003.jpg}<5>\includegraphics[width=5.5cm]{D:/home/jim/JamesAshe/talks/AMS_sectional_Knoxville_2014/thesis_pictures/cp_28501/cpShot-004.jpg}<6>\includegraphics[width=5.5cm]{D:/home/jim/JamesAshe/talks/AMS_sectional_Knoxville_2014/thesis_pictures/cp_28501/cpShot-005.jpg}<7>\includegraphics[width=5.5cm]{D:/home/jim/JamesAshe/talks/AMS_sectional_Knoxville_2014/thesis_pictures/cp_28501/cpShot-006.jpg}<8>\includegraphics[width=5.5cm]{D:/home/jim/JamesAshe/talks/AMS_sectional_Knoxville_2014/thesis_pictures/cp_28501/cpShot-007.jpg}<9>\includegraphics[width=5.5cm]{D:/home/jim/JamesAshe/talks/AMS_sectional_Knoxville_2014/thesis_pictures/cp_28501/cpShot-008.jpg}<10>\includegraphics[width=5.5cm]{D:/home/jim/JamesAshe/talks/AMS_sectional_Knoxville_2014/thesis_pictures/cp_28501/cpShot-009.jpg}<11>\includegraphics[width=5.5cm]{D:/home/jim/JamesAshe/talks/AMS_sectional_Knoxville_2014/thesis_pictures/cp_28501/cpShot-010.jpg}<12>\includegraphics[width=5.5cm]{D:/home/jim/JamesAshe/talks/AMS_sectional_Knoxville_2014/thesis_pictures/cp_28501/cpShot-011.jpg}<13>\includegraphics[width=5.5cm]{D:/home/jim/JamesAshe/talks/AMS_sectional_Knoxville_2014/thesis_pictures/cp_28501/cpShot-012.jpg}<14>\includegraphics[width=5.5cm]{D:/home/jim/JamesAshe/talks/AMS_sectional_Knoxville_2014/thesis_pictures/cp_28501/cpShot-014.jpg}<15>
\par\end{flushleft}%
\end{minipage}
\end{columns}
\end{frame}

\begin{frame}{}

All images created with ${\tt CirclePack}$

\href{http://www.math.utk.edu/~kens/CirclePack/}{\textcolor{blue}{http://www.math.utk.edu/$\sim$kens/CirclePack/}}

\begin{center}
\includegraphics[height=6.5cm]{pics/AMS_Knoxville_2014_pics/CP_screenshot_3_09}
\par\end{center}
\end{frame}

\end{document}
